\chapter{Resultados}
\label{chap:resultados}

Nesse capítulo será descrito o ambiente de execução dos experimentos, os resultados ob
tidos de cada algoritmo, bem como uma análise comparativa entre eles.

\section{Ambiente Computacional}

Os experimentos foram conduzidos em um ambiente computacional com as seguintes especificações 
técnicas, utilizando o \gls{WSL2} sobre o sistema operacional 
Windows 11:

\begin{itemize}
    \item \textbf{Processador:} AMD Ryzen 5 4600G com Radeon Graphics, arquitetura Zen 2, 6 núcleos 
    físicos com tecnologia \textit{Simultaneous Multithreading} (12 \textit{threads}), frequência 
    de 3.693 GHz;
    \item \textbf{Memória RAM:} 8 GiB, com 7.7 GiB disponíveis para aplicações;
    \item \textbf{Hierarquia de Cache:} L1 de 384 KiB, L2 de 3 MiB e L3 de 4 MiB;
    \item \textbf{Sistema Operacional:} Ubuntu 24.04.3 LTS (Noble Numbat) executado sobre 
    \textit{kernel} Linux 6.6.87.2-microsoft-standard-WSL2;
    \item \textbf{Compilador:} GCC versão 11.2.0;
    \item \textbf{Armazenamento:} Disco virtual de 1 TiB com sistema de arquivos ext4.
\end{itemize}

A escolha deste ambiente justifica-se pela disponibilidade de recursos computacionais adequados 
para a execução dos algoritmos propostos, bem como pela compatibilidade com ferramentas de 
desenvolvimento em Python. O uso do \gls{WSL2} permite executar nativamente aplicações Linux em ambiente 
Windows, combinando a flexibilidade do ecossistema Linux com a interface familiar do Windows.


\section{Execução dos Algoritmos}

